\documentclass[11pt,a4paper]{article}
\usepackage[utf8]{inputenc}
\usepackage[T1]{fontenc}
\usepackage[margin=1in]{geometry}
\usepackage{longtable}
\usepackage{booktabs}
\usepackage{enumitem}
\usepackage{hyperref}
\usepackage{xcolor}
\usepackage{array}
\usepackage{tabularx}

\hypersetup{colorlinks=true,linkcolor=black,citecolor=blue,urlcolor=blue}

\definecolor{red}{RGB}{220,38,38}
\definecolor{amber}{RGB}{217,119,6}
\definecolor{green}{RGB}{22,163,74}

\newcommand{\verdict}[1]{\textcolor{red}{\textbf{#1}}}
\newcommand{\warn}[1]{\textcolor{amber}{\textbf{#1}}}
\newcommand{\ok}[1]{\textcolor{green}{\textbf{#1}}}

\title{Lux Papers Corpus Audit\\[4pt]
\large 117 \texttt{.tex} files, 76{,}152 lines, evidence-based analysis}
\author{Scientist Agent}
\date{April 13, 2026}

\begin{document}
\maketitle
\tableofcontents
\newpage

%% =========================================================================
\section{Question}
%% =========================================================================

How consistent, non-redundant, and technically rigorous is the Lux papers
corpus? Specifically: (1) which papers are duplicates or near-duplicates,
(2) which contain AI-generated filler, (3) which deviate from the canonical
preamble, and (4) what is the recommended action for each?

%% =========================================================================
\section{Method}
%% =========================================================================

\begin{enumerate}
\item Read every \texttt{.tex} file's first 80 lines (title, author, date,
      documentclass, preamble, abstract).
\item Extract structural metadata: \texttt{documentclass} options, paper size,
      cover page invocation, bibliography style, \texttt{\textbackslash input\{shared/lstlang\}},
      column layout.
\item Grep for AI slop patterns: marketing language, LLM giveaway phrases,
      uncited superlatives.
\item Compare abstracts within each topical cluster to identify overlap.
\item Diff exact-duplicate candidates byte-for-byte.
\item Count \texttt{\textbackslash cite\{\}} and \texttt{\textbackslash bibitem}
      per file to identify uncited papers.
\end{enumerate}

All file paths are absolute under \texttt{/Users/z/work/lux/papers/}.

%% =========================================================================
\section{Inventory: 117 Papers by Topical Cluster}
%% =========================================================================

\subsection{Cluster 1: Core Consensus (7 papers)}
\begin{enumerate}[leftmargin=2em]
\item \texttt{lux-consensus.tex} --- Physics-inspired metastable consensus (Slush/Flare/Wave family). 1622 lines. Canonical.
\item \texttt{lux-quasar-consensus.tex} --- Quasar: quantum-secure triple-proof consensus. 701 lines.
\item \texttt{lux-quantum-consensus.tex} --- PQ secure multi-consensus with ML-DSA/ML-KEM/SLH-DSA (FIPS 205, formerly SPHINCS+). 557 lines.
\item \texttt{lp-020-quasar-consensus.tex} --- LP-020: Quasar formal specification (BLS + Pulsar + ML-DSA). 1398 lines. Canonical LP.
\item \texttt{lp-307-quantum-consensus-protocol.tex} --- LP-307: Eisenstein-series quantum consensus (research direction). 465 lines.
\item \texttt{lux-triple-proof-consensus.tex} --- Triple-proof quantum finality (BLS + STARK-ML-DSA + Pulsar). 323 lines.
\item \texttt{lux-fpc-consensus.tex} --- Fast Probabilistic Consensus with adaptive thresholds. 970 lines.
\end{enumerate}

\subsection{Cluster 2: Physics Sub-Protocols (8 papers)}
\begin{enumerate}[leftmargin=2em]
\item \texttt{lux-wave-protocol.tex} --- Decision stability / confidence monotonicity in DAG consensus. 712 lines.
\item \texttt{lux-photon-protocol.tex} --- Zero-latency pipelined binary consensus. 364 lines.
\item \texttt{lux-flare-protocol.tex} --- Certificate/skip conviction-based finality. 520 lines.
\item \texttt{lux-nova-protocol.tex} --- Supernova burst DAG consensus. 371 lines.
\item \texttt{lux-nebula-protocol.tex} --- Probabilistic finality for sparse networks. 409 lines.
\item \texttt{lux-prism-protocol.tex} --- Multi-spectrum parallel consensus. 432 lines.
\item \texttt{lux-ray-protocol.tex} --- Adaptive finality with light-speed confirmation. 405 lines.
\item \texttt{lux-field-protocol.tex} --- EM-field-inspired validator dipole consensus. 432 lines.
\end{enumerate}

\subsection{Cluster 3: MPC / Threshold Cryptography (6 papers)}
\begin{enumerate}[leftmargin=2em]
\item \texttt{lux-threshold-mpc.tex} --- Unified FROST + CGGMP21 for omnichain custody. 1368 lines. Canonical.
\item \texttt{lux-universal-threshold-signatures.tex} --- Universal threshold sigs (CMP + FROST + LSS + Doerner). 975 lines.
\item \texttt{lux-lss-mpc.tex} --- LSS: dynamic resharing framework. 332 lines.
\item \texttt{lux-mchain-mpc.tex} --- M-Chain: on-chain MPC custody coordination. 886 lines.
\item \texttt{lux-validator-mpc.tex} --- Threshold crypto for validator networks. 637 lines.
\item \texttt{mpc-custody.tex} --- Non-custodial securities custody via threshold signing. 518 lines.
\end{enumerate}

\subsection{Cluster 4: Teleport / Bridges (7 papers)}
\begin{enumerate}[leftmargin=2em]
\item \texttt{lux-teleport-architecture.tex} --- Engineering architecture, wire format, chain adapters. 786 lines. Canonical.
\item \texttt{lux-teleport-omnichain.tex} --- Sovereign omnichain liquidity protocol. 589 lines.
\item \texttt{lux-teleport-protocol.tex} --- LITP: light-client cross-chain transfers. 858 lines.
\item \texttt{teleport/main.tex} --- Teleport: omnichain bridge (April 2026 rewrite). 640 lines.
\item \texttt{lux-bridge.tex} --- Lux Bridge with MPC + Quasar hybrid verification. 1149 lines.
\item \texttt{lux-warp-messaging.tex} --- Intra-Lux BLS aggregate cross-appchain messaging. 805 lines.
\item \texttt{lux-secure-messaging.tex} --- Authenticated cross-chain message envelope protocol. 370 lines.
\end{enumerate}

\subsection{Cluster 5: FHE (6 papers)}
\begin{enumerate}[leftmargin=2em]
\item \texttt{lux-fhe-smart-contracts.tex} --- FHE VM architecture for private DeFi. 661 lines.
\item \texttt{lux-fhe-api.tex} --- Go API and developer model. 304 lines.
\item \texttt{lux-fhe-benchmarks.tex} --- Performance benchmarks and cost models. 314 lines.
\item \texttt{lux-fhe-mpc-hybrid.tex} --- FHE + MPC hybrid for trustless private computation. 330 lines.
\item \texttt{lux-tfhe.tex} --- TFHE: torus threshold FHE with distributed decryption. 475 lines.
\item \texttt{threshold-fhe-compliance.tex} --- Threshold FHE for regulatory forced decryption. 623 lines.
\end{enumerate}

\subsection{Cluster 6: Post-Quantum Cryptography (8 papers)}
\begin{enumerate}[leftmargin=2em]
\item \texttt{lux-quantum-threat-blockchain.tex} --- Systematic PQ threat model. 571 lines.
\item \texttt{lux-hybrid-pq-architecture.tex} --- Hybrid classical + PQ architecture. 553 lines.
\item \texttt{lux-pq-migration.tex} --- Three-phase PQ migration strategy. 349 lines.
\item \texttt{lux-crypto-agility.tex} --- Live algorithm rotation without forks. 230 lines.
\item \texttt{lux-corona-pq.tex} --- Pulsar: lattice-based threshold signatures. 480 lines.
\item \texttt{lux-ethfalcon-post-quantum.tex} --- ETHFALCON: EVM-optimized FALCON variants. 914 lines.
\item \texttt{lux-lattice-cryptography-primer.tex} --- Lattice crypto builder's guide. 619 lines.
\item \texttt{lux-hybrid-certificates.tex} --- Dual-sig X.509 for validators. 211 lines.
\end{enumerate}

\subsection{Cluster 7: EVM Execution / GPU (6 papers)}
\begin{enumerate}[leftmargin=2em]
\item \texttt{gpu-evm-whitepaper.tex} --- cevm: GPU-native EVM with PQ + FHE. 986 lines. Canonical.
\item \texttt{evmgpu-benchmark.tex} --- GPU-accelerated EVM benchmarks on M1 Max. 697 lines.
\item \texttt{dag-evm-formal.tex} --- DAG-EVM and conflict-graph consensus. 530 lines.
\item \texttt{lux-evm-precompiles.tex} --- Custom EVM precompile suite. 804 lines.
\item \texttt{lux-evm-precompile-engineering.tex} --- Precompile engineering tutorial. 583 lines.
\item \texttt{lux-pq-crypto-suite.tex} --- ML-KEM/ML-DSA/SLH-DSA as native precompiles. 248 lines.
\end{enumerate}

\subsection{Cluster 8: NTT (2 papers)}
\begin{enumerate}[leftmargin=2em]
\item \texttt{lux-ntt-transform.tex} --- NTT on EVM with 85\% gas reduction. 733 lines.
\item \texttt{lux-ntt-simd-implementation.tex} --- NTT SIMD implementation tutorial. 495 lines.
\end{enumerate}

\subsection{Cluster 9: Governance / DAO (2 papers)}
\begin{enumerate}[leftmargin=2em]
\item \texttt{lux-dao-governance-framework.tex} --- Azorius-based modular DAO framework. 1638 lines.
\item \texttt{lux-governance-dao.tex} --- Multi-ecosystem DAO with holographic consensus. 1100 lines.
\end{enumerate}

\subsection{Cluster 10: Identity (2 papers)}
\begin{enumerate}[leftmargin=2em]
\item \texttt{lux-id-did-specification.tex} --- DID specification (did:lux:address). 1092 lines.
\item \texttt{lux-id-iam.tex} --- Lux.id: full IAM system (Casdoor-based). 1138 lines.
\end{enumerate}

\subsection{Cluster 11: Economics (3 papers)}
\begin{enumerate}[leftmargin=2em]
\item \texttt{lux-tokenomics.tex} --- LUX token supply, rewards, fee-burning. 803 lines.
\item \texttt{lux-economics.tex} --- Appchain economics (validator attraction, fee revenue). 1011 lines.
\item \texttt{lux-validator-economics.tex} --- Validator incentives, slashing, delegation. 826 lines.
\end{enumerate}

\subsection{Cluster 12: DeFi / Markets (8 papers)}
\begin{enumerate}[leftmargin=2em]
\item \texttt{lux-lightspeed-dex.tex} --- High-frequency DEX. 683 lines.
\item \texttt{lux-perpetuals-derivatives.tex} --- Perpetuals with automated risk management. 866 lines.
\item \texttt{lux-omnichain-yield.tex} --- Yield-bearing bridge tokens. 943 lines.
\item \texttt{lux-liquid-staking.tex} --- Tokenized validator positions. 904 lines.
\item \texttt{lux-restaking.tex} --- Restaking for security extension. 1139 lines.
\item \texttt{lux-smart-order-routing.tex} --- Multi-provider execution routing. 676 lines.
\item \texttt{lux-privacy-pool.tex} --- Compliant privacy pools. 810 lines.
\item \texttt{lux-sovereign-defi.tex} --- Per-chain governance, Shariah compliance, ERC-3643. 965 lines.
\end{enumerate}

\subsection{Cluster 13: Local-First / Capsule / Web5 (5 papers)}
\begin{enumerate}[leftmargin=2em]
\item \texttt{lux-capsule-architecture.tex} --- Capsule: sovereign encrypted process. 1054 lines. Canonical.
\item \texttt{lux-app-architecture.tex} --- Identical to capsule-architecture (see \S\ref{sec:dupes}). 1053 lines.
\item \texttt{lux-firebase-exit.tex} --- Firebase Exit migration path. 794 lines.
\item \texttt{lux-web5-private-apps.tex} --- Web5 local-first architecture overview. 104 lines.
\item \texttt{lux-sovereign-keys.tex} --- Key management for local-first apps. 569 lines.
\end{enumerate}

\subsection{Cluster 14: Formal Verification / Testing (6 papers)}
\begin{enumerate}[leftmargin=2em]
\item \texttt{lux-formal-verification-lean4.tex} --- Lean 4 verification tutorial. 734 lines.
\item \texttt{lux-model-checking-tla.tex} --- TLA+ model checking tutorial. 577 lines.
\item \texttt{lux-protocol-verification-tamarin.tex} --- Tamarin security protocol verification. 531 lines.
\item \texttt{lux-symbolic-execution-halmos.tex} --- Halmos symbolic execution for contracts. 588 lines.
\item \texttt{lux-proof-methodology.tex} --- Proof obligations and verification tiers. 285 lines.
\item \texttt{lux-fuzz-testing-methodology.tex} --- Property-based testing and fuzzing. 484 lines.
\end{enumerate}

\subsection{Cluster 15: Chain Architecture (5 papers)}
\begin{enumerate}[leftmargin=2em]
\item \texttt{lux-whitepaper.tex} --- The Lux Network whitepaper. 535 lines. Canonical.
\item \texttt{lux-chain-architecture.tex} --- 12-chain architecture overview. 204 lines.
\item \texttt{lux-zchain.tex} --- Z-Chain: ZK-UTXO privacy chain. 493 lines.
\item \texttt{lux-gchain-graphql.tex} --- G-Chain: GraphQL query engine. 753 lines.
\item \texttt{lux-achain-attestation.tex} --- A-Chain: TEE attestation for AI compute. 494 lines.
\end{enumerate}

\subsection{Cluster 16: Infrastructure / LP Standards (5 papers)}
\begin{enumerate}[leftmargin=2em]
\item \texttt{lp-102-encrypted-sqlite-replication.tex} --- Encrypted SQLite replication. 1075 lines.
\item \texttt{lp-103-omnichain-explorer.tex} --- Single-binary block explorer. 309 lines.
\item \texttt{lp-104-graph-engine.tex} --- Embedded subgraph-compatible GraphQL. 247 lines.
\item \texttt{lp-checkpoint-anchor.tex} --- On-chain checkpoint anchoring. 430 lines.
\item \texttt{chain-first-architecture.tex} --- Blockchain as source of truth for regulated finance. 531 lines.
\end{enumerate}

\subsection{Cluster 17: Quantum LP Series (5 papers)}
\begin{enumerate}[leftmargin=2em]
\item \texttt{lp-305-quantum-secure-assets.tex} --- LP-305: Quantum secure assets. 449 lines.
\item \texttt{lp-306-lx-social-dex.tex} --- LP-306: LX social trading exchange. 302 lines.
\item \texttt{lp-308-quantum-amm.tex} --- LP-308: QuantumSwap AMM. 509 lines.
\item \texttt{lp-309-quantum-network-comparison.tex} --- LP-309: Quantum network comparison. 422 lines.
\end{enumerate}

\subsection{Cluster 18: Security / Operations (7 papers)}
\begin{enumerate}[leftmargin=2em]
\item \texttt{lux-master-security-model.tex} --- Comprehensive security model. 669 lines.
\item \texttt{lux-cross-chain-security.tex} --- Multi-chain security architecture. 353 lines.
\item \texttt{lux-hsm-boundary.tex} --- HSM and threshold boundary design. 221 lines.
\item \texttt{lux-zero-trust-validators.tex} --- Zero-trust validator operations. 229 lines.
\item \texttt{lux-reproducible-builds.tex} --- Signed software and build pipelines. 262 lines.
\item \texttt{lux-performance-security-tradeoffs.tex} --- Performance vs security tradeoffs. 416 lines.
\item \texttt{lux-smart-contract-auditing.tex} --- Systematic auditing methods. 582 lines.
\end{enumerate}

\subsection{Cluster 19: Data / State (4 papers)}
\begin{enumerate}[leftmargin=2em]
\item \texttt{lux-data-availability.tex} --- Data availability sampling. 1105 lines.
\item \texttt{lux-state-sync.tex} --- Merkle range proof state sync. 801 lines.
\item \texttt{lux-verkle-trees.tex} --- Verkle trees for stateless clients. 918 lines.
\item \texttt{lux-fraud-proofs.tex} --- Optimistic rollup fraud proofs. 1327 lines.
\end{enumerate}

\subsection{Cluster 20: Miscellaneous (8 papers)}
\begin{enumerate}[leftmargin=2em]
\item \texttt{ethos-of-lux.tex} --- Eight design principles. 285 lines.
\item \texttt{lux-adoption-roadmap.tex} --- Research-to-deployment roadmap. 289 lines.
\item \texttt{lux-agent-sovereignty.tex} --- AI agents as first-class processes. 564 lines.
\item \texttt{lux-agentic-consensus-zap.tex} --- Multi-agent AI consensus over ZAP. 267 lines.
\item \texttt{lux-decentralized-cloud.tex} --- Providers as relays/storage/compute markets. 861 lines.
\item \texttt{lux-market-nft.tex} --- NFT marketplace aggregator. 754 lines.
\item \texttt{lux-oracle-infrastructure.tex} --- Decentralized price feeds. 823 lines.
\item \texttt{lux-tee-computing-mesh.tex} --- TEE confidential computing infrastructure. 675 lines.
\end{enumerate}

\subsection{Non-Papers (2 files)}
\begin{enumerate}[leftmargin=2em]
\item \texttt{cover-lux.tex} --- Standalone cover page. 36 lines. Not a paper.
\item \texttt{lux-threshold-ux.tex} --- Threshold UX product design. 812 lines.
\item \texttt{lux-trust-kernel.tex} --- Trust kernel: identity + key policy + threshold governance. 849 lines.
\item \texttt{lux-distributed-systems-testing.tex} --- Testing from unit to Jepsen. 669 lines.
\item \texttt{lux-security-token-standard.tex} --- ERC-1400 on PQ EVM with MPC custody. 676 lines.
\end{enumerate}

\textbf{Total: 117 \texttt{.tex} files} (115 papers + 1 shared file + 1 standalone cover).

%% =========================================================================
\section{Duplicates and Overlaps}\label{sec:dupes}
%% =========================================================================

\subsection{EXACT DUPLICATE: lux-app-architecture.tex = lux-capsule-architecture.tex}

\verdict{These are the same paper.} Both have identical title (``The Capsule:
A Sovereign Encrypted Process for Web5 Computing''), identical abstract,
identical author (Zach Kelling), identical date (2025). A byte-level diff
shows only 2 differences: (1) \texttt{capsule-architecture} adds one sentence
with a \texttt{\textbackslash cite\{lux-tokenomics\}} reference, and (2) it
replaces UTF-8 tree-drawing characters (\texttt{\char"2514, \char"251C,
\char"2502}) with ASCII equivalents (\texttt{|--, +--}).

\textbf{Action:} Delete \texttt{lux-app-architecture.tex}. Keep
\texttt{lux-capsule-architecture.tex} (it has the citation fix).

\subsection{NEAR-DUPLICATE: teleport/main.tex $\approx$ lux-teleport-omnichain.tex}

Both cover the same protocol: Lux Teleport non-custodial omnichain bridge.
Abstracts are nearly identical (``270+ blockchains,'' ``FROST and CGGMP21,''
same three innovations). Section structures differ by 2 sections.
\texttt{teleport/main.tex} is the April 2026 rewrite (640 lines, twocolumn,
different preamble), \texttt{lux-teleport-omnichain.tex} is the 2022
(revised December 2025) version (589 lines, onecolumn).

\textbf{Action:} Keep \texttt{teleport/main.tex} as canonical (newest). Delete
\texttt{lux-teleport-omnichain.tex} or archive it as v1.

\subsection{HIGH OVERLAP: Consensus Cluster (7 papers, 3 redundant)}

The consensus cluster has significant redundancy:

\begin{itemize}
\item \texttt{lux-consensus.tex} (1622 lines) --- the comprehensive canonical
      paper covering the full Slush/Flare/Wave/Nova/Nebula/Prism/Quasar stack
      with Lean 4 proofs.
\item \texttt{lp-020-quasar-consensus.tex} (1398 lines) --- the formal LP
      specification of Quasar. Legitimate companion to the above.
\item \texttt{lux-fpc-consensus.tex} (970 lines) --- FPC as the ``wave'' layer.
      Focused sub-protocol paper. Legitimate.
\item \verdict{lux-quasar-consensus.tex} (701 lines) --- an earlier, less
      formal treatment of Quasar. Superseded by LP-020. Claims ``4,761,904 TPS
      with finality'' without reproducible benchmark. \textbf{Superseded.}
\item \verdict{lux-quantum-consensus.tex} (557 lines) --- early 2022 paper
      claiming ``first high-performance blockchain ready for the quantum era''
      and ``50,000+ TPS.'' Superseded by LP-020 and LP-307. Contains Keywords
      section (unusual for the corpus). \textbf{Superseded.}
\item \texttt{lp-307-quantum-consensus-protocol.tex} (465 lines) ---
      Eisenstein-series quantum consensus. Explicitly labels itself a ``2023
      research direction'' with status section. Legitimate exploratory paper.
\item \texttt{lux-triple-proof-consensus.tex} (323 lines) --- focused on the
      BLS + STARK-ML-DSA + Pulsar triple-proof mechanism. Uses non-standard
      preamble (see \S\ref{sec:structure}). Could be merged into LP-020.
\end{itemize}

\subsection{HIGH OVERLAP: MPC Cluster (6 papers, 2 redundant)}

\begin{itemize}
\item \texttt{lux-threshold-mpc.tex} (1368 lines) --- canonical unified
      framework paper covering FROST + CGGMP21 + chain adapters.
\item \texttt{lux-universal-threshold-signatures.tex} (975 lines) --- covers
      CMP + FROST + LSS + Doerner. Broader scope but overlaps 60--70\% with
      threshold-mpc. Claims ``over four years of continuous development
      (February 2021 - August 2025, 800+ commits)'' which is marketing.
\item \texttt{lux-lss-mpc.tex} (332 lines) --- focused on LSS dynamic
      resharing. Legitimate sub-topic paper.
\item \texttt{lux-mchain-mpc.tex} (886 lines) --- M-Chain: the on-chain
      coordination layer for MPC. Distinct enough from the crypto papers.
\item \verdict{lux-validator-mpc.tex} (637 lines) --- threshold crypto for
      validators. Heavily overlaps with threshold-mpc. Covers the same CGGMP21
      + FROST + DKG content but scoped to validators. \textbf{Merge into
      threshold-mpc.}
\item \texttt{mpc-custody.tex} (518 lines) --- non-custodial securities
      custody. Distinct application focus (2-of-3 threshold for user wallets).
\end{itemize}

\subsection{HIGH OVERLAP: Teleport Cluster (4+ papers, 2 redundant)}

Beyond the near-duplicate noted above:
\begin{itemize}
\item \texttt{lux-teleport-architecture.tex} (786 lines) --- engineering spec:
      wire format, chain adapters, multi-tenant. Canonical.
\item \texttt{lux-teleport-protocol.tex} (858 lines) --- LITP: light-client
      verification, Warp intra-ecosystem, lock-and-mint. Legitimate protocol
      layer paper.
\item \texttt{lux-bridge.tex} (1149 lines) --- ``Lux Bridge'' with MPC +
      Quasar hybrid verification. The bibliography explicitly notes it is
      superseded by \texttt{lux-teleport-omnichain.tex}. \textbf{Superseded.}
\item \texttt{lux-warp-messaging.tex} (805 lines) --- BLS aggregate
      cross-appchain messaging. Distinct protocol (intra-Lux only).
\item \texttt{lux-secure-messaging.tex} (370 lines) --- message envelope
      specification. Could be an appendix to teleport-architecture.
\end{itemize}

\subsection{HIGH OVERLAP: DAO Cluster (2 papers)}

\begin{itemize}
\item \texttt{lux-dao-governance-framework.tex} (1638 lines) --- Azorius
      framework, ERC-4337, Hats Protocol. Technical implementation paper.
\item \texttt{lux-governance-dao.tex} (1100 lines) --- multi-ecosystem DAO
      with holographic consensus. Higher-level governance design. Uses manual
      numbered references instead of BibTeX. Contains specific claims
      (``94 LPs successfully activated,'' ``\$128M treasury,'' ``12,847 active
      participants'') without external citation.
\end{itemize}

Both are DAO governance papers but at different abstraction levels. The
governance-dao paper is significantly more bloated (36 slop-word hits).
\textbf{Merge governance-dao's unique content into dao-governance-framework.}

\subsection{MODERATE OVERLAP: EVM Precompiles (3 papers)}

\begin{itemize}
\item \texttt{lux-evm-precompiles.tex} (804 lines) --- full precompile suite
      (BLS12-381, BN254, batch ecrecover, PQ, Warp).
\item \texttt{lux-evm-precompile-engineering.tex} (583 lines) --- engineering
      tutorial on building precompiles. Distinct pedagogical purpose.
\item \texttt{lux-pq-crypto-suite.tex} (248 lines) --- ML-KEM/ML-DSA/SLH-DSA
      precompiles specifically. Subset of evm-precompiles but more focused on
      the three FIPS standards.
\end{itemize}

Minor overlap. All three can coexist if pq-crypto-suite is positioned as the
focused PQ-only companion to the broader precompiles paper.

\subsection{MODERATE OVERLAP: Local-First (3+ papers)}

\texttt{lux-capsule-architecture.tex}, \texttt{lux-firebase-exit.tex},
\texttt{lux-web5-private-apps.tex}, and \texttt{lux-sovereign-keys.tex} all
cover the local-first encrypted-SQLite architecture. They are at different
abstraction levels: web5 (104 lines) is a 2-page overview, firebase-exit is the
migration narrative, capsule is the full spec, sovereign-keys is the key
management layer. Acceptable stratification but web5-private-apps is too thin to
justify a standalone paper.

\subsection{ACCEPTABLE: FHE Cluster (6 papers)}

Six papers covering distinct sub-topics: API, benchmarks, MPC hybrid, smart
contracts, TFHE, compliance. This is a well-stratified cluster with clear
separation of concerns.

\subsection{ACCEPTABLE: Physics Sub-Protocols (8 papers)}

Each of Photon, Wave, Flare, Nova, Nebula, Prism, Ray, Field covers a distinct
consensus sub-protocol with different algorithmic contributions. The main
overlap risk is with \texttt{lux-consensus.tex} which surveys all of them, but
these are deeper treatments of individual protocols.

%% =========================================================================
\section{AI Slop / Hyperbole}
%% =========================================================================

\subsection{Worst Offenders (slop-word density per 1000 lines)}

\begin{tabular}{lrrp{6cm}}
\toprule
\textbf{File} & \textbf{Lines} & \textbf{Slop hits} & \textbf{Worst phrases} \\
\midrule
\texttt{lux-governance-dao.tex} & 1100 & 36 & ``holistic,'' ``comprehensive,'' ``seamless,'' ``ecosystem'' (26 uses), ``robust foundation,'' ``blueprint for multi-ecosystem governance'' \\
\texttt{lux-dao-governance-framework.tex} & 1638 & 30 & ``cutting-edge technologies,'' ``solid foundation for the next generation,'' ``significant advancement'' \\
\texttt{lux-perpetuals-derivatives.tex} & 866 & 20 & ``comprehensive,'' ``seamless,'' ``innovative,'' ``leveraged'' (sometimes legitimate financial term) \\
\texttt{lux-id-iam.tex} & 1138 & 19 & ``unprecedented challenges,'' ``comprehensive,'' ``seamless integration'' \\
\texttt{lux-market-nft.tex} & 754 & 17 & ``seamless,'' ``comprehensive,'' ``next-generation,'' ``cornerstone of the evolving NFT economy,'' ``positioning Lux.market as\ldots'' \\
\texttt{lux-liquid-staking.tex} & 904 & 15 & ``comprehensive,'' ``innovative,'' ``seamless,'' ``ecosystem'' (14 uses) \\
\texttt{lux-oracle-infrastructure.tex} & 823 & 13 & ``robust foundation for the next generation of decentralized applications, bridging the gap between on-chain logic and real-world data'' \\
\texttt{lux-fraud-proofs.tex} & 1327 & 12 & ``positions Lux at the forefront,'' ``unprecedented opportunities,'' ``robust fraud proof framework'' \\
\bottomrule
\end{tabular}

\subsection{Specific Offending Phrases (with file:line)}

\textbf{Marketing superlatives:}
\begin{itemize}[leftmargin=2em, itemsep=0pt]
\item ``positions Lux Network at the forefront of blockchain security for the next generation'' --- \texttt{lux-quasar-consensus.tex:603}
\item ``unprecedented challenges for identity'' --- \texttt{lux-id-iam.tex:61}
\item ``unprecedented opportunities for decentralized AI services'' --- \texttt{lux-fraud-proofs.tex:1186}
\item ``unprecedented composability'' --- \texttt{lux-gchain-graphql.tex:657}
\item ``cutting-edge technologies'' --- \texttt{lux-dao-governance-framework.tex:1542}
\item ``cornerstone of the evolving NFT economy'' --- \texttt{lux-market-nft.tex:715}
\item ``blueprint for the next generation of universal identity platforms'' --- \texttt{lux-id-did-specification.tex:760}
\item ``bridging the gap between cryptographic theory and blockchain practicality'' --- \texttt{lux-ethfalcon-post-quantum.tex} abstract
\item ``paving the way for quantum-resistant blockchain systems'' --- \texttt{lux-ntt-transform.tex} abstract
\item ``This work represents the first production-ready NTT implementation for EVM'' --- \texttt{lux-ntt-transform.tex} abstract
\item ``solid foundation for the next generation of decentralized organizations'' --- \texttt{lux-dao-governance-framework.tex:1556}
\end{itemize}

\textbf{LLM giveaway patterns:}
\begin{itemize}[leftmargin=2em, itemsep=0pt]
\item No instances of ``In the rapidly evolving landscape'' --- \ok{clean}
\item No instances of ``delving deeper'' --- \ok{clean}
\item 2 instances of ``Additionally'' (legitimate usage) --- \ok{acceptable}
\item Extensive use of ``comprehensive'' (219 total hits for tracked slop words) concentrated in 8 papers
\end{itemize}

\textbf{Uncited specific claims (in governance-dao):}
\begin{itemize}[leftmargin=2em, itemsep=0pt]
\item ``94 LPs successfully activated'' --- no source
\item ``\$128M treasury managed transparently'' --- no source
\item ``12,847 active participants'' --- no source
\item ``18.3\% voter participation (4-5x higher than typical DAOs)'' --- no source for the comparison
\end{itemize}

%% =========================================================================
\section{Structural Inconsistencies}\label{sec:structure}
%% =========================================================================

\subsection{Preamble / Cover Page}

\begin{tabular}{lrp{6cm}}
\toprule
\textbf{Pattern} & \textbf{Count} & \textbf{Notes} \\
\midrule
Uses \texttt{shared/luxcover} + \texttt{\textbackslash luxcoverpage} & 108 & Canonical pattern \\
Uses \texttt{\textbackslash maketitle} instead of luxcover & 7 & lp-103, lp-104, lux-chain-architecture, lux-web5-private-apps, lux-sovereign-defi, lux-omnichain-yield, teleport/main.tex \\
No cover at all & 2 & lux-agentic-consensus-zap, lux-triple-proof-consensus \\
Standalone cover file & 1 & cover-lux.tex (not referenced by any paper) \\
\bottomrule
\end{tabular}

\texttt{lux-agentic-consensus-zap.tex} and \texttt{lux-triple-proof-consensus.tex}
use a completely different preamble: \texttt{fancyhdr} with \texttt{\textbackslash rhead/\textbackslash lhead},
\texttt{titlesec} custom section formatting, \texttt{parskip}, manual
\texttt{\textbackslash begin\{center\}} title block. No \texttt{shared/luxcover},
no \texttt{\textbackslash title\{\}}, no \texttt{\textbackslash author\{\}}, no
\texttt{\textbackslash date\{\}}.

\subsection{Paper Size}

\begin{tabular}{lr}
\toprule
\textbf{Paper size} & \textbf{Count} \\
\midrule
\texttt{a4paper} & 81 \\
\texttt{letterpaper} & 6 (lightspeed-dex, pq-crypto-suite, quantum-consensus, tee-computing-mesh, zap-benchmarks, zap-wire-protocol) \\
No size specified & 25 \\
\bottomrule
\end{tabular}

\subsection{Column Layout}

\begin{tabular}{lr}
\toprule
\textbf{Layout} & \textbf{Count} \\
\midrule
Single column (default) & 110 \\
\texttt{twocolumn} & 5 (lux-web5-private-apps, teleport/main.tex, lux-omnichain-yield, lux-threshold-mpc, lux-universal-threshold-signatures, lux-sovereign-defi) \\
\bottomrule
\end{tabular}

\subsection{Bibliography Style}

\begin{tabular}{lrp{6cm}}
\toprule
\textbf{Style} & \textbf{Count} & \textbf{Notes} \\
\midrule
\texttt{bibliographystyle\{plain\}} & 32 & Most common \\
\texttt{bibliographystyle\{plainnat\}} & 4 & consensus, economics, restaking, web5 \\
\texttt{thebibliography} only (inline) & 77 & \\
Manual numbered list (\texttt{\textbackslash item}) & 1 & lux-governance-dao.tex \\
No bibliography at all & 4 & lux-zchain, cover-lux, lux-fraud-proofs, lux-governance-dao \\
\bottomrule
\end{tabular}

\subsection{Date Formats}

Five distinct date formats in use:
\begin{enumerate}[leftmargin=2em]
\item Year only: ``2024'' (most common, 60+ papers)
\item Month Year: ``February 2025'' (12 papers)
\item Original/Revised with \texttt{\textbackslash quad|\textbackslash quad}: ``Original: 2022\quad|\quad Revised: February 2026'' (12 papers)
\item Year with parenthetical revision: ``2022 (Revised December 2025)'' (3 papers)
\item Empty: 2 papers (teleport/main.tex, lux-agentic-consensus-zap.tex)
\end{enumerate}

\subsection{Author Attribution}

\begin{tabular}{lrp{6cm}}
\toprule
\textbf{Author format} & \textbf{Count} & \textbf{Notes} \\
\midrule
Zach Kelling (sole, with thanks/affiliation) & $\sim$90 & Standard \\
Zach Kelling + Vishnu Seesahai & 12 & bridge, teleport-*, threshold-*, sovereign-defi, etc. \\
``Lux Architecture Team'' & 2 & lp-103, lp-104 \\
``Lux Industries, Inc.'' (in \texttt{\textbackslash begin\{center\}} block) & 2 & lux-agentic-consensus-zap, lux-triple-proof-consensus \\
Vishnu J. Seesahai (sole) & 1 & lux-lss-mpc \\
No \texttt{\textbackslash author} command & 2 & agentic-consensus-zap, triple-proof-consensus \\
\bottomrule
\end{tabular}

\subsection{Papers Without \texttt{\textbackslash input\{shared/lstlang\}}}

37 papers that include code listings lack \texttt{\textbackslash input\{shared/lstlang\}}.
This is not always an error (some define their own minimal lstset), but papers
using Solidity, Go, Rust, or TypeScript code should use the shared definitions.

\subsection{Papers With Zero \texttt{\textbackslash cite\{\}} Commands}

37 of 115 papers (32\%) contain zero inline citations. They have
\texttt{\textbackslash bibitem} entries in a bibliography section but never
reference them with \texttt{\textbackslash cite\{\}} in the text body. This
means every bibliographic reference is decorative --- readers cannot trace any
claim to a specific source. Notable offenders:

\begin{itemize}[leftmargin=2em, itemsep=0pt]
\item \texttt{chain-first-architecture.tex} (531 lines, 0 cites)
\item \texttt{lux-fraud-proofs.tex} (1327 lines, 0 cites)
\item \texttt{lux-governance-dao.tex} (1100 lines, 0 cites)
\item \texttt{lux-id-iam.tex} (1138 lines, 0 cites --- has bibitems but never cites them)
\item \texttt{lux-gchain-graphql.tex} (753 lines, 0 cites)
\item \texttt{lux-oracle-infrastructure.tex} (823 lines, 0 cites)
\item \texttt{lux-perpetuals-derivatives.tex} (866 lines, 0 cites)
\end{itemize}

%% =========================================================================
\section{Recommendations}
%% =========================================================================

\subsection{DELETE (5 files)}

\begin{enumerate}
\item \texttt{lux-app-architecture.tex} --- exact duplicate of capsule-architecture.
\item \texttt{lux-teleport-omnichain.tex} --- superseded by teleport/main.tex.
\item \texttt{lux-bridge.tex} --- self-identifies as superseded in its own bibliography.
\item \texttt{cover-lux.tex} --- standalone cover page, not referenced by any paper.
\item \texttt{lux-quasar-consensus.tex} --- superseded by lp-020-quasar-consensus.tex. Contains unsubstantiated ``4.7M TPS'' claim.
\end{enumerate}

\subsection{MERGE (4 actions)}

\begin{enumerate}
\item Merge \texttt{lux-quantum-consensus.tex} into \texttt{lp-020-quasar-consensus.tex} --- the LP is the canonical Quasar spec; the 2022 paper adds nothing the LP does not cover.
\item Merge \texttt{lux-validator-mpc.tex} into \texttt{lux-threshold-mpc.tex} --- validator-specific threshold content belongs in the unified framework paper.
\item Merge \texttt{lux-governance-dao.tex} into \texttt{lux-dao-governance-framework.tex} --- the DAO-framework paper is more technical; the governance-dao paper's unique content (holographic consensus, LP process, cross-ecosystem coordination) should be absorbed as sections.
\item Merge \texttt{lux-secure-messaging.tex} into \texttt{lux-teleport-architecture.tex} --- the message envelope spec is part of the Teleport architecture.
\end{enumerate}

\subsection{REWRITE (8 papers need significant cleanup)}

\begin{enumerate}
\item \texttt{lux-market-nft.tex} --- marketing document masquerading as a paper. 17 slop-word hits. ``cornerstone of the evolving NFT economy.'' Rewrite as concise technical spec or delete.
\item \texttt{lux-oracle-infrastructure.tex} --- conclusion is pure marketing (``robust foundation for the next generation''). Strip marketing language, add inline citations.
\item \texttt{lux-fraud-proofs.tex} --- 1327 lines with zero \texttt{\textbackslash cite\{\}}. ``unprecedented opportunities.'' ``positions Lux at the forefront.'' Add citations, remove superlatives.
\item \texttt{lux-id-iam.tex} --- ``unprecedented challenges.'' 14 uses of ``ecosystem.'' 0 inline citations despite 10 bibitems. Wire up citations.
\item \texttt{lux-gchain-graphql.tex} --- ``unprecedented composability.'' ``next-generation multi-chain experiences.'' 0 inline citations. Strip conclusion slop.
\item \texttt{lux-ntt-transform.tex} --- ``paving the way'' and ``first production-ready'' in abstract. Remove superlatives. The technical content (85\% gas reduction) stands on its own.
\item \texttt{lux-perpetuals-derivatives.tex} --- 0 inline citations, 20 slop-word hits. Heavy on ``comprehensive'' and ``seamless.''
\item \texttt{lux-liquid-staking.tex} --- 15 slop-word hits, 14 uses of ``ecosystem.'' Tighten language.
\end{enumerate}

\subsection{STANDARDIZE PREAMBLE (2 papers need immediate fix)}

\texttt{lux-agentic-consensus-zap.tex} and \texttt{lux-triple-proof-consensus.tex}
use a completely non-standard preamble (fancyhdr, titlesec, manual center-block
title, no \texttt{\textbackslash title/author/date}, no luxcover). Convert to
the canonical template.

\subsection{Canonical Preamble Template}

All Lux papers should use this structure:

\begin{verbatim}
\documentclass[11pt,a4paper]{article}
\usepackage{shared/luxcover}
\usepackage[utf8]{inputenc}
\usepackage[T1]{fontenc}
\usepackage{amsmath,amssymb,amsfonts,amsthm}
\usepackage{hyperref}
\usepackage{booktabs}
\usepackage{listings}
\usepackage{xcolor}
\input{shared/lstlang}    % when code listings are used
% ... additional packages as needed ...

\title{Paper Title}
\author{
  Zach Kelling\thanks{zach@lux.network}\\
  \textit{Lux Industries}\\
  \texttt{research@lux.network}
}
\date{Original: YYYY\quad|\quad Revised: Month YYYY}

\begin{document}
\luxcoverpage
\begin{abstract} ... \end{abstract}
\tableofcontents
\newpage
\end{verbatim}

Key rules:
\begin{enumerate}
\item Always \texttt{a4paper}, never \texttt{letterpaper} (6 papers to fix).
\item Always onecolumn unless the paper is genuinely short (5 papers to assess).
\item Always \texttt{\textbackslash luxcoverpage}, never \texttt{\textbackslash maketitle} (7 papers to fix).
\item Always use \texttt{\textbackslash cite\{\}} to reference bibitems (37 papers to fix).
\item Date format: ``Original: Month YYYY\texttt{\textbackslash quad|\textbackslash quad} Revised: Month YYYY'' for revised papers, ``Month YYYY'' for new papers.
\item Author: ``Zach Kelling'' (with co-authors where applicable), never ``Lux Architecture Team.''
\end{enumerate}

\subsection{Priority Actions}

\begin{enumerate}
\item \textbf{Immediate:} Delete the 5 files listed above.
\item \textbf{Week 1:} Fix the 2 non-standard preambles, 7 maketitle-instead-of-luxcover papers, and 6 letterpaper papers.
\item \textbf{Week 2:} Execute the 4 merges.
\item \textbf{Week 3:} Rewrite the 8 papers flagged for slop.
\item \textbf{Ongoing:} Wire up \texttt{\textbackslash cite\{\}} in the 37 papers that have bibitems but zero inline citations.
\end{enumerate}

%% =========================================================================
\section{Sources}
%% =========================================================================

All findings are based on direct file inspection of \texttt{/Users/z/work/lux/papers/}.
Line numbers cited are from the files as of April 13, 2026. No external sources
were consulted; this is a structural/content audit, not a correctness review.

\end{document}
